\documentclass[10pt,a4paper]{amsart}
\usepackage[margin=19mm]{geometry}
\usepackage{amsmath,amssymb,mathtools}
\usepackage[scr=rsfs,cal=euler]{mathalfa}
\usepackage{xfrac}
\usepackage[unicode,hidelinks,bookmarks=false]{hyperref}
\newcommand{\ie}{i.e.\ }
\newcommand{\cf}{cf.\ }
\newcommand{\resp}{respectively\ }
\newcommand{\Subsection}{\subsection}
\providecommand{\nobreakdash}{\nobreak\hbox{-}}
\setlength{\emergencystretch}{2em}
\multlinegap0pt
\usepackage{amssymb}
\usepackage{enumerate}
\usepackage{xcolor}
\usepackage{tikz}
\newtheorem{lemma}{Lemma}
\newtheorem{remark}{Remark}
\newcommand{\Ac}{\mathcal{A}}
\title{Study of Rota-Baxter Operators in Matrix $C^*$-Algebras Motivated by Toeplitz Structures, and Applications to Sliding Mode Control}
\author{Marwa Ennaceur}
\date{Source: arXiv:2605.08126v1 (CC BY 4.0)}
\begin{document}
\maketitle

\noindent\emph{Source and licence.} Study of Rota-Baxter Operators in Matrix $C^*$-Algebras Motivated by Toeplitz Structures, and Applications to Sliding Mode Control. Version arXiv:2605.08126v1 (CC BY 4.0), distributed by arXiv under the Creative Commons Attribution 4.0 International licence, http://creativecommons.org/licenses/by/4.0/, as recorded in the arXiv licence metadata for this exact version and shown on the abstract page https://arxiv.org/abs/2605.08126v1. Full licence notice: \texttt{licenses/arxiv-2605.08126-CC-BY-4.0.txt}.\par\smallskip

\noindent\textit{Editorial scope.} Counted unit: the article's lemma lem:lie (source0.tex lines 451--458). The article's complete original proof of it is delivered with it (source0.tex lines 459--669, one \texttt{proof} environment of 211 lines). No mathematics is written, repaired or re-derived: the statement and the proof are the article's own lines, in the article's own notation, and no retained line is substituted or re-flowed.  Retained uncounted as the source-local context the proof needs: lines 167--169.  Nothing was omitted from any retained span, so the package carries no omission record.  The retained lines cite 1 works, whose entries are transcribed verbatim from the article's own bibliography source in the e-print and delivered with short editorial labels recorded in the item's reference mapping.  No retained line contains a figure environment, an image-inclusion command or drawing code, so no image is delivered and the package declares no asset dependency.  Editorial formatting only, in addition: an AMS \texttt{amsart} preamble that restates the article's own declarations and macros used by the retained lines, namely the environments \texttt{lemma}, \texttt{remark} on separate counters, together with the standard packages those lines need.  The retained environments print in this document as Lemma 1, Remark 1 in order of appearance, whereas the article prints the same items with the numbering of its own full document.  The complete native source of the article as distributed in the arXiv e-print for this exact version is bundled unchanged as \texttt{sources/200138/source0.tex}.
\begin{equation}\label{eq:RB}
P(x)P(y) = P(x P(y)) + P(P(x) y) + \lambda P(x y), \quad \forall\, x,y\in\Ac.
\end{equation}

\begin{lemma}[Induced Lie Structure]\label{lem:lie}
Let \(P\) be a Rota-Baxter operator of weight \(\lambda\) on an associative algebra
\(\Ac\). Define the bracket
\begin{equation}\label{eq:bracket}
[x,y]_P := [P(x), y] + [x, P(y)] + \lambda [x,y].
\end{equation}
Then \((\Ac, [\cdot,\cdot]_P)\) is a Lie algebra.
\end{lemma}

\begin{proof}
\textbf{Antisymmetry.} Using antisymmetry of the ordinary commutator,
\[
[y,x]_P = [P(y),x]+[y,P(x)]+\lambda[y,x]
         = -[x,P(y)]-[P(x),y]-\lambda[x,y]
         = -[x,y]_P.
\]

\textbf{Jacobi identity.} We verify
\[
[[x,y]_P,z]_P + [[y,z]_P,x]_P + [[z,x]_P,y]_P = 0.
\]
The key algebraic ingredient is obtained by writing the Rota-Baxter identity
\eqref{eq:RB} for the pair $(x,y)$ and for the pair $(y,x)$:
\begin{align*}
P(x)P(y) &= P(xP(y)) + P(P(x)y) + \lambda P(xy),\\
P(y)P(x) &= P(yP(x)) + P(P(y)x) + \lambda P(yx).
\end{align*}
Subtracting the second from the first and using $[a,b] = ab - ba$:
\begin{align*}
[P(x),P(y)]
  &= P(xP(y)) - P(yP(x)) + P(P(x)y) - P(P(y)x)
     + \lambda P(xy) - \lambda P(yx)\\
  &= P([x,P(y)]) + P([P(x),y]) + \lambda P([x,y]).
\end{align*}
Rearranging yields the identity
\begin{equation}\label{eq:RB_comm}
P\bigl([P(x),y]\bigr) + P\bigl([x,P(y)]\bigr)
  = \bigl[P(x),P(y)\bigr] - \lambda P\bigl([x,y]\bigr),
\end{equation}
which we use below. Set $u := [x,y]_P = [P(x),y]+[x,P(y)]+\lambda[x,y]$.

\textbf{Step 1: Compute $P(u)$.}
By linearity of $P$ and identity \eqref{eq:RB_comm}:
\begin{align*}
P(u)
  &= P([P(x),y]) + P([x,P(y)]) + \lambda P([x,y]) \\
  &= \bigl([P(x),P(y)] - \lambda P([x,y])\bigr) + \lambda P([x,y]) \\
  &= [P(x),P(y)].
\end{align*}
Thus $P([x,y]_P) = [P(x),P(y)]$, i.e.\ $P$ maps the $[\cdot,\cdot]_P$-bracket to the
ordinary commutator of the $P$-images.

\textbf{Step 2: Expand $[[x,y]_P,z]_P$.}
By definition \eqref{eq:bracket} with $u=[x,y]_P$:
\begin{align*}
[[x,y]_P,z]_P
  &= [P(u),z] + [u,P(z)] + \lambda[u,z] \\
  &= [[P(x),P(y)],z]
     + \bigl[[P(x),y]+[x,P(y)]+\lambda[x,y],\; P(z)\bigr]
     + \lambda\bigl[[P(x),y]+[x,P(y)]+\lambda[x,y],\; z\bigr].
\end{align*}

\textbf{Step 3: Cyclic sum.}
The cyclic sum is taken over the three even permutations
$(x,y,z) \to (y,z,x) \to (z,x,y)$, which we denote $\sum_{\mathrm{cyc}(x,y,z)}$.
Taking this cyclic sum and expanding, we obtain five families of terms:
\begin{align*}
J &:= \sum_{\mathrm{cyc}}[[x,y]_P,z]_P \\
  &= \underbrace{\sum_{\mathrm{cyc}}[[P(x),P(y)],z]}_{S_1}
   + \underbrace{\sum_{\mathrm{cyc}}[[P(x),y],P(z)]
                +\sum_{\mathrm{cyc}}[[x,P(y)],P(z)]}_{S_2}
   + \underbrace{\lambda\sum_{\mathrm{cyc}}[[x,y],P(z)]}_{G^{(2)}}\\
  &\quad + \underbrace{\lambda\sum_{\mathrm{cyc}}[[P(x),y],z]}_{G^{(1)}}
   + \underbrace{\lambda\sum_{\mathrm{cyc}}[[x,P(y)],z]}_{H}
   + \underbrace{\lambda^2\sum_{\mathrm{cyc}}[[x,y],z]}_{G^{(3)}}.
\end{align*}

\textbf{Step 4: $S_1 + S_2 = 0$.}
For any three elements $A,B,C$ in an associative algebra,
\begin{equation}\label{eq:jacobi_triple}
[[A,B],C] + [[B,C],A] + [[C,A],B] = 0.
\end{equation}
The six terms of $S_2$ are:
\begin{align}\label{eq:group2}
S_2 &= \underbrace{[[P(x),y],P(z)]}_{T_1}
      + \underbrace{[[P(y),z],P(x)]}_{T_2}
      + \underbrace{[[P(z),x],P(y)]}_{T_3} \notag\\
    &\quad + \underbrace{[[x,P(y)],P(z)]}_{T_4}
      + \underbrace{[[y,P(z)],P(x)]}_{T_5}
      + \underbrace{[[z,P(x)],P(y)]}_{T_6}.
\end{align}
Apply \eqref{eq:jacobi_triple} to $(P(x),P(y),z)$:
$[[P(x),P(y)],z] = -T_2-T_6$.\\
Apply \eqref{eq:jacobi_triple} to $(P(y),P(z),x)$:
$[[P(y),P(z)],x] = -T_3-T_4$.\\
Apply \eqref{eq:jacobi_triple} to $(P(z),P(x),y)$:
$[[P(z),P(x)],y] = -T_1-T_5$.\\
Summing: $S_1 = -(T_1+T_2+T_3+T_4+T_5+T_6) = -S_2$, so $S_1+S_2 = 0$.

\begin{remark}\label{rem:S1S2}
Neither $S_1$ nor $S_2$ vanishes individually in general; their \emph{joint}
cancellation $S_1+S_2=0$ is what matters.
\end{remark}

\textbf{Step 5: $G^{(1)} + H + G^{(2)} = 0$.}
Apply the ordinary Jacobi identity \eqref{eq:jacobi_triple} to the triple
$(P(x)+\lambda x,\, y,\, z)$:
\[
  [[P(x)+\lambda x,y],z]+[[y,z],P(x)+\lambda x]+[[z,P(x)+\lambda x],y]=0.
\]
Expanding linearly and taking the cyclic sum $\sum_{\mathrm{cyc}(x,y,z)}$ over the
three even permutations $(x,y,z)\to(y,z,x)\to(z,x,y)$:
\begin{align*}
  \sum_{\mathrm{cyc}}\bigl([[P(x),y],z]+\lambda[[x,y],z]\bigr)
  +\sum_{\mathrm{cyc}}\bigl([[y,z],P(x)]+\lambda[[y,z],x]\bigr)
  +\sum_{\mathrm{cyc}}\bigl([[z,P(x)],y]+\lambda[[z,x],y]\bigr)=0.
\end{align*}
Each $\lambda$-coefficient cyclic sum vanishes by \eqref{eq:jacobi_triple}
(applied to $(x,y,z)$). Thus:
\begin{equation}\label{eq:group3_combined}
  \underbrace{\sum_{\mathrm{cyc}}[[P(x),y],z]}_{G^{(1)}/\lambda}
  + \underbrace{\sum_{\mathrm{cyc}}[[y,z],P(x)]}_{(*)}
  + \underbrace{\sum_{\mathrm{cyc}}[[z,P(x)],y]}_{(**)} = 0.
\end{equation}
We now identify the sums $(*)$ and $(**)$ with $G^{(2)}/\lambda$ and $H/\lambda$
respectively. For $(*)$: writing out the three terms of the cyclic sum over
$(x,y,z)\to(y,z,x)\to(z,x,y)$ gives
\[
  (*) = [[y,z],P(x)] + [[z,x],P(y)] + [[x,y],P(z)].
\]
Under the relabelling $(x,y,z)\mapsto(z,x,y)$ (the second even permutation applied
once more), the generic summand $[[x,y],P(z)]$ maps to $[[z,x],P(y)]$, and cycling
again gives $[[y,z],P(x)]$. Hence the three terms of $(*)$ coincide with the three
terms of $\sum_{\mathrm{cyc}}[[x,y],P(z)] = G^{(2)}/\lambda$ (same set of terms,
in a different order). Thus $(*)=G^{(2)}/\lambda$.

For $(**)$: writing out the cyclic sum gives
\[
  (**) = [[z,P(x)],y]+[[x,P(y)],z]+[[y,P(z)],x].
\]
Relabelling $(x,y,z)\mapsto(y,z,x)$ sends the generic summand $[[x,P(y)],z]$ to
$[[y,P(z)],x]$, and cycling again gives $[[z,P(x)],y]$. Hence the three terms of
$(**)$ coincide with the three terms of $\sum_{\mathrm{cyc}}[[x,P(y)],z] = H/\lambda$.
Thus $(**)=H/\lambda$.

Multiplying \eqref{eq:group3_combined} by $\lambda$ gives $G^{(1)}+H+G^{(2)}=0$.

\begin{remark}[Correct cancellation in Group 3 --- explicit verification]
\label{rem:jacobi_correction}
Neither $G^{(1)}$, $H$, nor $G^{(2)}$ vanishes individually in general.
We verify this explicitly for $P = P_+$ ($\lambda = -1$) in $M_2(\mathbb{C})$
with $x = e_{12}$, $y = e_{21}$, $z = e_{11}$.

Recall $P_+(e_{12}) = e_{12}$ (upper triangular), $P_+(e_{21}) = 0$ (strictly lower),
$P_+(e_{11}) = e_{11}$ (diagonal, hence upper triangular).

\medskip
\noindent\textbf{Computation of $G^{(1)}$} $= \lambda\sum_{\mathrm{cyc}}[[P(x),y],z]$
$= \lambda\bigl([[P(e_{12}),e_{21}],e_{11}] + [[P(e_{21}),e_{11}],e_{12}]
+ [[P(e_{11}),e_{12}],e_{21}]\bigr)$.

\begin{itemize}
  \item $[[e_{12},e_{21}],e_{11}] = [e_{11}-e_{22},e_{11}] = 0$ (diagonal matrices commute).
  \item $[[0,e_{11}],e_{12}] = 0$.
  \item $[e_{11},e_{12}] = e_{12}$, so $[[e_{11},e_{12}],e_{21}] = [e_{12},e_{21}]
        = e_{11}-e_{22}$.
\end{itemize}
Thus $G^{(1)}/\lambda = e_{11}-e_{22}$, giving $G^{(1)} = (-1)(e_{11}-e_{22}) = e_{22}-e_{11}$.

\medskip
\noindent\textbf{Computation of $G^{(2)}$} $= \lambda\sum_{\mathrm{cyc}}[[x,y],P(z)]$
$= \lambda\bigl([[e_{12},e_{21}],P(e_{11})] + [[e_{21},e_{11}],P(e_{12})]
+ [[e_{11},e_{12}],P(e_{21})]\bigr)$.

\begin{itemize}
  \item $[[e_{12},e_{21}],e_{11}] = [e_{11}-e_{22},e_{11}] = 0$.
  \item $[e_{21},e_{11}] = e_{21}$, so $[[e_{21},e_{11}],e_{12}] = [e_{21},e_{12}]
        = e_{22}-e_{11}$.
  \item $[[e_{11},e_{12}],0] = 0$.
\end{itemize}
Thus $G^{(2)}/\lambda = e_{22}-e_{11}$, giving $G^{(2)} = (-1)(e_{22}-e_{11}) = e_{11}-e_{22}$.

\medskip
\noindent\textbf{Computation of $H$} $= \lambda\sum_{\mathrm{cyc}}[[x,P(y)],z]$
$= \lambda\bigl([[e_{12},0],e_{11}] + [[e_{21},e_{11}],e_{12}]
+ [[e_{11},e_{12}],e_{21}]\bigr)$.

\begin{itemize}
  \item $[[e_{12},0],e_{11}] = 0$.
  \item $[[e_{21},e_{11}],e_{12}] = [e_{21},e_{12}] = e_{22}-e_{11}$.
  \item $[e_{11},e_{12}] = e_{12}$, so $[[e_{11},e_{12}],e_{21}] = [e_{12},e_{21}]
        = e_{11}-e_{22}$.
\end{itemize}
Thus $H/\lambda = (e_{22}-e_{11})+(e_{11}-e_{22}) = 0$, giving $H = 0$.

\medskip
\noindent\textbf{Final verification:}
\[
G^{(1)}+H+G^{(2)} = (e_{22}-e_{11}) + 0 + (e_{11}-e_{22}) = 0. \quad\checkmark
\]

This confirms: (i) $G^{(1)}/\lambda = e_{11}-e_{22} \neq 0$,
(ii) $G^{(2)}/\lambda = e_{22}-e_{11} \neq 0$,
(iii) $H = 0$ for this triple, and (iv) their joint sum $G^{(1)}+H+G^{(2)}=0$.
The identity is a \emph{joint} cancellation, not three separate vanishings.
A fortiori, the invalid argument
$\sum_{\mathrm{cyc}}[[A,B],L(C)] = L(\sum_{\mathrm{cyc}}[[A,B],C])$ (which would
require $L$ to act on the full expression, not only on $C$) cannot be used to simplify
these sums.
\end{remark}

\textbf{Step 6: $G^{(3)}=0$.}
$G^{(3)} = \lambda^2\bigl([[x,y],z]+[[y,z],x]+[[z,x],y]\bigr) = 0$
by \eqref{eq:jacobi_triple} applied to $(A,B,C)=(x,y,z)$.

\medskip
\textbf{Conclusion.} Since $S_1+S_2=0$, $G^{(1)}+H+G^{(2)}=0$, and $G^{(3)}=0$,
the total cyclic sum $J=0$, establishing the Jacobi identity for $[\cdot,\cdot]_P$.
See also \cite[Theorem~3.3]{BaiGuoNi2011}.
\end{proof}

\begin{thebibliography}{99}
\bibitem[BaiGuo]{BaiGuoNi2011} C.~Bai, L.~Guo, and X.~Ni, Lie algebras associated with Rota-Baxter operators, \textit{Commun. Contemp. Math.} \textbf{13} (2011), no.~3, 461--490.
\end{thebibliography}
\end{document}
